\section{Lean mechanization and implementation boundary}\label{sec:lean}
The artifact uses Lean 4.14.0 with Mathlib pinned by a checked-in Lake
manifest (Mathlib commit \texttt{4bbdccd9c5f8}). The library is a
mechanization of finite operational states, certificate conditions and
an asynchronous event system. It is not a formalization of a SELL
sequent calculus, a verified Python compiler or a proof of the measured
Python program's execution time.

\paragraph{Representations.}
Bits are \texttt{ZMod 2}, coordinates are finite types, and an
observation is a function into optional bits. The formal
\texttt{PeelingState} stores the partial assignment $K$; its unknown
neighbors, degrees and syndromes are derived functions. This gives
coherence by construction. The Python implementation instead maintains
cached arrays and adjacency lists. The Lean update lemmas prove the
mathematical degree and syndrome equations, but there is no verified
refinement of those Python caches.

A formal asynchronous context stores $K$, finite pending sets $Q_j$
and an ordered event observer. Its residual fields are derived from
$(K,Q)$, and its well-formedness condition is
$\Nb(j)\cap U\subseteq Q_j\subseteq\Nb(j)$. Thus exact multiplicities
of the canonical tokens are enforced by representation. The relation
to the raw multisets and their exact multiplicity invariant is proved
in Section~\ref{sec:sem}; the Lean development operates on the
canonical representation rather than mechanizing a quotient of all
malformed raw multisets. The log field records chronological events;
its multiset image represents the manuscript's unordered log atoms.
A raw context alone does not determine that order. The correspondence
therefore relates a raw context $M$ together with a declared history
$T$, whose multiset image is the log of $M$, to the canonical context
$(K,Q,T)$; it does not reconstruct $T$ from the final multiset.

\begin{table}[t]
\centering\small
\caption{Representation and evidence boundaries. Arrows between the
layers require the correspondence stated here; they are not all
mechanized refinements.}
\label{tab:representation-boundary}
\begin{tabular}{>{\raggedright\arraybackslash}p{.21\textwidth}>{\raggedright\arraybackslash}p{.31\textwidth}>{\raggedright\arraybackslash}p{.39\textwidth}}
\toprule
Layer & Representation & Evidence for the correspondence \\
\midrule
Raw token execution & Admissible multiset $M$ and declared history $T$ &
Exact multiplicities, delayed residual fields and unordered log:
manuscript invariant and proof in Section~\ref{sec:sem}. \\
Canonical asynchronous Lean & Partial assignment $K$, pending sets $Q_j$,
ordered observer $T$ &
Invariant preservation, arbitrary-interleaving soundness and exact
trace completeness: kernel-checked for this representation. \\
Canonical operational Lean & State $K$ with derived syndromes and degrees;
typed traces and witnesses &
Kernel-checked replay relation and success/failure certificate theorems. \\
Python checker and adapter & Cached arrays, adjacency lists and serialized
residual witnesses &
Finite differential agreement with the typed Lean checker; no general
proof of parsing, cache refinement or program equivalence. \\
\bottomrule
\end{tabular}
\end{table}

The relation \texttt{RewStar H C T D} has the event list $T$ as an
explicit index. Its soundness theorem accepts arbitrary interleavings
of Propagate and Solve from a well-formed context, including endpoints
with unpropagated known incidences. Completeness constructs a derivation
for the very same prescribed $T$, ending at
\texttt{encodeLogged} of the folded operational state with the previous
log followed by $T$. It does not assert equality with an encoding that
discards the Solve events, and it does not attempt to extract ordered
data by eliminating a proof-irrelevant reachability proposition.
The abstract adequacy principle in Section~\ref{sec:sem} is a manuscript
result. The Lean library proves the corresponding concrete LDPC
properties; it contains no generic adequacy theorem for arbitrary
labelled transition systems.

\begin{table}[t]
\centering\small
\caption{Mechanized results and their principal declarations.
The accompanying verification report records full statements,
build commands and axiom-dependency output.}
\label{tab:lean}
\begin{tabular}{>{\raggedright\arraybackslash}p{.37\textwidth}>{\raggedright\arraybackslash}p{.54\textwidth}}
\toprule
Property & Lean declaration(s) \\
\midrule
Forced assignments preserve all completions &
\texttt{forced\_value}, \texttt{trace\_extends} \\
Replay matches the inductive trace relation &
\texttt{replayFull\_eq\_some\_iff} \\
Left-kernel separator soundness and existence &
\texttt{lem\_unsat\_witness\_sound},
\texttt{lem\_unsat\_witness\_complete} \\
Accepted success and failure certificates &
\texttt{thm\_sound\_success}, \texttt{thm\_sound\_failure} \\
Residual and empty-trace certificate completeness &
\texttt{residual\_unsat\_iff\_witness},
\texttt{thm\_complete\_failure}, \texttt{thm\_complete\_success} \\
Real-valued BEC likelihood maximization &
\texttt{bec\_completion\_maximizes}, \texttt{thm\_ml} \\
Invariant and local operational correspondence &
\texttt{propagate\_wf}, \texttt{solve\_wf}, \texttt{solve\_valid} \\
Propagation normalization and order independence &
\texttt{propagation\_normalization},
\texttt{propagation\_confluence}, \texttt{propagate\_commute} \\
Arbitrary-interleaving trace soundness &
\texttt{rewstar\_sound}, \texttt{adequacy\_soundness} \\
Logged trace completeness and same-trace equivalence &
\texttt{adequacy\_completeness}, \texttt{thm\_adequacy\_complete} \\
\bottomrule
\end{tabular}
\end{table}

\paragraph{Algebra and likelihood.}
The separator-existence proof uses finite-dimensional linear duality:
a vector outside the image of the residual matrix is separated by a
linear functional, represented by a row vector over $\Ftwo$. This is
an existence theorem, not a formal verification of the Python row
eliminator or its transformation matrix. The residual system has a
mechanized solution-or-separator equivalence. The named checker
completeness theorems construct accepted certificates using the empty
trace; the stronger manuscript statement for every prescribed valid
partial trace follows from its residual-bijection proof and is not
presented as a single checker-completeness theorem in Lean.
Success and failure soundness
are proved for the executable typed checker; failure soundness uses
preservation of every compatible completion by the replayed trace.

The likelihood function is real-valued. It is defined as zero on an
incompatible word, and otherwise as
$\varepsilon^{e}(1-\varepsilon)^k$, where $e$ counts erasures and
$k=n-e$ counts visible entries.
Lean proves that an accepted completion maximizes this function over
codewords when $0\leq\varepsilon\leq1$. This is the closed form of the
coordinate-product model in Proposition~\ref{prop:core-ml}; the
coordinate-product factorization itself is justified there rather than
by a separate Lean probability-space construction. The endpoint values
are included without asserting a conditional distribution on a
zero-probability observation. Neither ML theorem has the vacuous
conclusion \texttt{True}.

\paragraph{What the formal checker guarantees.}
The success witness in Lean is a total function on variable indices;
only its residual entries are used. The Python certificate instead
serializes those entries in increasing residual-index order. The
adapter pads known coordinates when translating a Python witness to
the typed formal input. The formal checker receives already typed
matrix, observation, trace and witness data. Consequently parsing,
binary-domain rejection, serialization and conversion to those types
form a separate implementation boundary. Finite differential cases
exercise that boundary and compare the Lean and Python acceptance
results. All 575 supplied differential cases agreed, including returned
words for accepted successes (369 acceptances and 206 rejections). The
run used Lean's interpreter; the native executable built on this host
could not be loaded by macOS 27. These finite comparisons do not prove
equivalence for all inputs.

The full verification report distinguishes source inspection,
compilation, theorem statements and their transitive axiom
dependencies. No project axiom or admitted proof is used to establish
the listed results. Lean's ordinary foundational assumptions and its
kernel, together with the pinned libraries, remain part of the formal
trust base. Execution of a compiled checker additionally relies on the
compiler and runtime. The raw-token representation correspondence,
SELL-label erasure, global rewrite termination including Solve events,
stopping-set and nullity characterizations,
Gaussian-elimination implementation, sparse complexity and bit-size
bounds, and experimental/statistical claims are not presented as
mechanized results. Their mathematical proofs or empirical evidence
are given separately.
